\section{Introduction}

The RLCE cryptosystem of Wang~\cite{DBLP:journals/corr/Wang15aa, RLCE-spec} was one of the
69 schemes submitted to NIST's postquantum cryptography standardization
process. RLCE builds its public key from a generalized Reed--Solomon code
$GRS_{n,k}$ by injecting $w$ \emph{twin pairs}: for each pair, two GRS columns
are mixed with two random columns via a $2 \times 2$ invertible matrix, and all
$N = n + w$ columns are permuted. The three unbroken parameter crucially all
have $w = n-k$, the exact boundary where the CLT square distinguisher provably
fails.

This scheme was one of many~\cite{DBLP:conf/pqcrypt/Wieschebrink09,
      DBLP:journals/dcc/BaldiBC13, DBLP:conf/wcc/Wieschebrink16,
      DBLP:journals/iet-it/AlmeidaBNS24, DBLP:conf/redundancy/Krouk20} attempts to
patch Niederreiter's original attempt to instantiate the
McEliece~\cite{McEliece1978} cryptosystem with GRS
codes~\cite{DBLP:journals/dc/Niederreiter86}, which was itself broken by the
Sidenikov--Shestakov attack~\cite{sidelnikov1992insecurity}. The original
scheme and all of these other previous attempts have fallen to cryptanalysis.

Due to security concerns raised by the Couvreur--Lequesne--Tillich
(CLT)~\cite{CLT18} attack on smaller parameters, NIST did not advance RLCE to
the second round; it appears among the schemes for which NIST ``lacked full
confidence in security''~\cite{Moody19}. See~\cite{NIST8240} for the Round~1
report. No subsequent public attack has been advanced against the $w = n-k$
parameter sets since, and due to the large parameter sizes the scheme may not
have received any attention at all in the intervening period. However, with the
recent improvements in AI cryptanalysis ~\cite{cryptoeprint:2026/1593,
      cryptoeprint:2026/2266} causing concerns from eminent cryptographers that key
exchange/key encapsulation itself may be at risk ~\cite{green2026pkc}, it is
likely important to subject even large parameter schemes to renewed
cryptanalysis attempts in order to enable a potential broadening of
standardized schemes for cryptographic agility as a hedge against AI
breakthroughs against lattices, isogenies or quasi-cyclic codes.

There has been recent hubbub over the Ghoshal--Ishai--Jain--Sun~\cite{GIJS26}
provable distinguisher and heuristic key recovery for the venerable and ancient
(by cryptographic standards) McEliece scheme~\cite{McEliece1978} and subsequent
improvement papers~\cite{cryptoeprint:2026/1984}. Most notably it has led to
the German Federal Office for Information Security (BSI) shifting its guidance
from wholesale approval to recommending against the use of the Classic McEliece
scheme in new deployments~\cite{bsi2026classicmceliece}. In the course of a
very long thread on NIST's PQC forum over exactly how devastating or not these
attacks are, in practice, Dr. Wang commented that "I do not think these attacks
on McEliece will work on RLCE. I still believe RLCE has sufficient randomness
in the generator matrix and is a good alternative for McEliece
schemes,"~\cite{wang2026pqcforum} which piqued the interest of this work's
author.

We contradict Wang's claim for \emph{both} distinguishing and key recovery by
using the GIJS distinguisher: the distinguishing step is proved, and the
key-recovery step is an experimentally verified cryptanalysis~\cite{GIJS26}.

\subsection{The Attack}\label{sec:intro-attack}

In slightly more detail, our attack proceeds as follows.
\begin{enumerate}
      \item \textbf{Shorten} the public key at $\ell = k - k'$ random positions,
            for
            some \emph{small} $k'$.
      \item \textbf{Distinguish} using the recent
            Ghoshal--Ishai--Jain--Sun~\cite{GIJS26}
            distinguisher (Theorem~\ref{thm:ghoshal}), which operates beyond the
            Schur-square barrier, applied \emph{directly} with every column as
            hold-out.
      \item \textbf{Group} columns across trials via co-occurrence with twin pairs.
      \item \textbf{Recover} the original GRS columns from 3 twin pairs.
      \item \textbf{Solve} using Sidelnikov-Shestakov.
\end{enumerate}

A key insight is that GIJS will still work even with some random columns, and
that, furthermore, the pairs of noisy columns provide useful information to the
GIJS distinguisher. We show in Section~\ref{sec:gijs} that the GIJS
distinguisher for GRS subcodes should work on shortened RLCE public keys
whenever $sC + U > dD$, where $s$ is the jet parameter, $d$ is the degree of the
multivariate polynomials, $C$ is the number of clean (either from shortening at
one of a twin pair's columns or from an initially clean column) columns in the
shortened code, and $U$ is the number of twin pairs remaining in the shortened
code.

Extensive experimental testing shows that the GRS distinguisher on the
parameter sets of interest essentially never fails to diagnose a ``clean" GRS
column as GRS and a noisy column as ``random." This leads to practical attacks.

In short, we simply choose uniformly at random columns to shorten at for a
relatively small (less than 100 for the sets considered) number of separate
\emph{trials}. In each trial, we run the GIJS distinguisher on the shortened
code, with each column held out separately as the holdout column. This enables
us to learn which columns in the shortening remain part of a twin pair (meaning
both it and its pair column are in the shortened code) in that shortening, and
refer to these columns as \emph{flagged}.

To recover twin pairs, we observe that (assuming no errors in the
distinguisher) a column is \emph{flagged} in a trial if and only its twin pair
column is also \emph{flagged} in that trial. This fact means that after a
relatively small number of randomized trials, for a given column A, there will
be only other column B that is flagged in and only in trials in which column A
is flagged, and that this column B is the twin pair for column A.

After recovering all twin pairs, we use some clever ideas in\ref{sec:fastp2} to
unmix 3 of these twin pairs back into actual GRS columns. We note that a
simpler route to do so would involve simply trying all $(q-1)^3$ non-degenerate
linear combinations of the pairs, one of which will correctly unmix all 3.

Finally, after we have recovered a $k \times (k+3)$ pure GRS matrix, we are able to
run the celebrated~\cite{sidelnikov1992insecurity} to recover the code
structure (the scrambling matrix $S$). At this point, unmixing the remainder of the twin pairs
becomes trivial (see~\cite{CLT18}). We end up with the full private key.

\subsection{Experimental Success}\label{sec:exp-success}

We were able to run a full key recovery (see Section~\ref{sec:total-costs}) at
the 128 bit parameters (ID-0) in the specification that took approximately 30.5
hours of core work on a standard high-end laptop to recover the key. We have
also given key recovery data for several other smaller parameter sets for
comparison in the $w=n-k$ regime; see Section~\ref{sec:total-costs} for
details.

While we did not execute an experimental attack on $ID-1$ and $ID-2$ due to a
lack of desire to spend personal funds, we estimate that a cloud based attack should be possible
for no more than a few hundred dollars; the biggest issues, as mentioned
in~\cite{GIJS26}, are the memory costs of the attack. However, these are easily
solvable in the cloud for the given parameter sets.

\subsection{Further Work}\label{sec:further-work}
The scheme is completely broken in the $w=n-k$ parameter regime, as a back of the envelope 
computation estimates that a public key would need to be in the petabytes with hundreds of megabytes per 
ciphertext in order to reach $128$ bits of security. However, there is a natural
defensive option; push the number of twin pairs $w$ to its maximum
$w=n$, eliminating pure columns entirely ($N=2n$). Our attack does not work in
this regime when $2k \geq n$, so one could potentially propose this as a patch.
However, due to the matrix becoming wider while the error correction regime
does not change, ISD attack perform better, necessitating a bump in $k$, and
significantly raising parameter sizes.

Below, we have calculated parameter sets that seem to be immune to the attacks
given in this paper, as well as to ISD-based attacks. 

\begin{table}[!ht]
      \caption{Proposed $w=n$ parameter sets immune to our attack and ISD}
      \begin{tabular}{lrrrrrrrrr}
            \toprule
                      & $n$  & $k$  & $t$ & $N$  & field    & Stern & pk (KB)   & ct (kB)\\
            \midrule
            128-fixed & 1191 & 619  & 286 & 2382 & $2^{11}$ & 140.2 & 1980    & 3.2 \\
            192-fixed & 1774 & 922  & 426 & 3548 & $2^{11}$ & 204.0 & 4393    & 4.8 \\
            256-fixed & 2355 & 1225 & 565 & 4710 & $2^{12}$ & 268.1 & 8452    & 6.9 \\
            \bottomrule
      \end{tabular}
\end{table}

We do not claim that these parameters are secure, and if forced, we would bet
against their security given the long, long history of attacks on GRS-McEliece
schemes and various attempts to patch. We merely give them as new targets and
the smallest scheme parameters secure given our new attacks for future research. 